\section{结论}
\label{sec:8}

我们提出了动态可组合性的形式化基础，将经典的效应与协效应概念提升为运行时机制。可逆效应处理局部时间可组合性：每个上下文变换携带一个运行时追踪的逆操作，且追踪和恢复均保持组合性，因此在组件移除时上下文得以恢复。响应式协效应处理局部空间可组合性：每当上下文变化时，组件根据其协效应规约收到通知，每次变更被分类为激活、去激活或中性，协效应隔离改变已声明键解析到的内容，协效应拦截改变绑定被使用的方式。我们将效应上下文和协效应上下文统一为单一的上下文类型，其中协效应上的观察等价性为效应提供了独立性，构成了时空可组合性的编程范式。将这些机制组合为组件的概念后，便得到动态组合的演算，其元理论将时空可组合性从单个组件推广到交织组件的整个系统。我们将该范式实现为 Cordis 元框架，其核心库提供效应追踪和协效应解析，声明式组件加载器提供配置调和与热模块替换。Koishi 案例研究在拥有超过 4000 个社区插件的生产系统中验证了 Cordis 的设计。

除了人工策划的插件生态系统之外，未来验证的一个引人注目的方向是自演化智能体框架（\S\ref{sec:1.2.2}），其中 AI 智能体持续地在极少人工监督下生成并替换其自身的框架组件。在该场景中应用 Cordis 将验证在快速组件替换下完全恢复的时间保证，以及在频繁拓扑变化下依赖协调的空间保证。这样的验证将展示该范式作为智能体框架和其他自主系统中可恢复、可协调且持续自演化基础的适用性。
